\begin{abstract}
Cross-silo federated learning (FL) between hospitals is limited by the uplink: every round is a scheduled transfer over an institutional WAN, a cellular slice or a satellite link. We treat the job as a network tenant with a communication quota and measure, over \num{TotalRuns} bit-reproducible runs, how the rounds in which a site transmits set accuracy, network cost and the client-level differential-privacy (DP) budget. \tierfed{} solves exactly the integer programme that minimises a local-SGD drift bound under the quota and keeps a per-site head private. At $\rho=0.25$ it reaches \num{diabetesmlpTierfedAuroc}~AUROC against \num{diabetesmlpFedperAuroc} for full-rate FedPer, \ifflag{diabetesmlpTierfedNonInferiorAuroc}{non-inferior}{not non-inferior} at a 0.01 margin with \num{diabetesmlpTierfedUploadSavingPct}\% less upload, and \num{compTierfedStcUploadMb}~MB per site with sparse ternary compression. Under DP the transmission schedule is itself observable, so schedulers that react to client updates, ours included, leak information the accountant does not charge. \tierfed{}-DP plans the period and noise multiplier before training from the privacy target, sampling rate and byte quota. At matched budgets $\varepsilon\le8$ with Poisson sampling, however, client-level DP reduces every method to near-chance AUROC ($\le$\num{dpmQHalfEpsEightFixedAuroc}); fewer releases help, a fixed period is best, and the planner with its default prior is not better than FedPer. On four real heart-disease institutions \tierfed{} matches full-rate training at a third of the upload, and we check the network model against paced-TCP emulation of a \num{wideParams}-parameter model (agreement within \num{emuMaxAbsRelErrPct}\%). A fixed period matches or exceeds \tierfed{} everywhere, tier differentiation is rarely active, attacks do not separate the methods, and an earlier negative claim is retracted.
\end{abstract}
